\begin{titlepage}
\newgeometry{margin=0.85in}
\thispagestyle{empty}

\begin{center}
{\large University of Rochester\\}
{\normalsize Warner School of Education\\[0.35in]}

{\color{PrimaryRed}\rule{\textwidth}{1.4pt}}\\[0.22in]

{\Huge\bfseries Indicator 13 Checklist Review\\[0.08in]
Transition Planning, Health Access,\\[0.08in]
and Student Agency}\\[0.18in]

{\Large Review of a Sample High School IEP}\\[0.22in]

{\color{PrimaryRed}\rule{0.78\textwidth}{0.7pt}}\\[0.32in]

\colorbox{SoftBlue}{%
\parbox{0.88\textwidth}{%
\centering\vspace{0.12in}
\textbf{From Compliance to Future-Ready Support}\\[0.05in]
A strengths-based review of transition planning for a student with complex health, communication, academic, and self-advocacy needs.
\vspace{0.12in}
}%
}
\end{center}

\vspace{0.38in}

\begin{center}
\renewcommand{\arraystretch}{1.18}
\setlength{\tabcolsep}{4pt}
\begin{tabular}{>{\bfseries\RaggedLeft}p{1.65in} p{4.9in}}
Student & Piter Zacari Garcia Bautista \\
Course & ED452C \\
Assignment & Indicator 13 Checklist Review \\
Focus & Transition planning for neurodivergent, chronically ill, disabled, trauma-experienced, and mental-health vulnerable students \\
IEP Reviewed & Sample High School IEP: Jimmy, age 17 \\
Submission Date & \today \\
\end{tabular}
\end{center}

\vspace{0.34in}

\begin{center}
\begin{minipage}{0.86\textwidth}
\itshape
This review evaluates the transition section of a sample high school IEP using Indicator 13 checklist expectations. The analysis considers what is present, the quality of the information included, what additional information would strengthen the plan, and whether the transition section meaningfully supports the student's postsecondary goals.

\vspace{0.18in}

The review is grounded in a course lens centered on access, agency, self-determination, belonging, chronic illness, neurodivergence, depression-responsive planning, trauma-informed support, and intersecting student identities. The goal is not only to ask whether the IEP meets compliance expectations, but whether it prepares the student for adult life with dignity and realistic support.
\end{minipage}
\end{center}

\restoregeometry
\end{titlepage}
\clearpage
\onehalfspacing

\section*{Introduction}

This review applies the Indicator 13 transition checklist to the course-provided Sample High School IEP for Jimmy, a 17-year-old student in Grade 10. The assignment requires a completed checklist and a reflection addressing what is present in the IEP, the quality of the transition information, and what additional information would strengthen the transition section's alignment with the rest of the IEP. It also asks whether the transition section is designed appropriately to meet the student's postsecondary goals and objectives.

Jimmy's IEP contains unusually detailed present-level information. He is identified as having Other Health Impairment and Speech Language Impairment; he has Prader-Willi syndrome, significant food-safety and temperature-regulation needs, language delays, academic support needs, social goals, counseling services, and a stated career interest connected to early childhood education. The strongest parts of the IEP protect access and safety. The weaker parts are where protection needs to become future-ready agency.

My analysis uses a strengths-based, neurodiversity-affirming, trauma-informed, and chronic-illness-aware lens. I do not read the sample as evidence that Jimmy has depression or a trauma history; the IEP does not document that. However, transition planning for chronically ill and neurodivergent students should anticipate stress, fatigue, social vulnerability, medical complexity, mental-health risk, and the emotional weight of adult-life transitions rather than waiting until those needs become crises.

\begin{figure}[H]
\centering
\begin{tikzpicture}[scale=0.95, every node/.style={font=\small, align=center}]
  \draw[rounded corners=6pt, very thick, PrimaryRed] (-4.9,-1.55) rectangle (4.9,1.75);

  \node[font=\small\bfseries, align=center, text width=8.7cm] at (0,1.35)
  {From compliance to student agency, health access, and adult-life support};

  \node[circle, fill=SoftBlue,  draw=DarkGray, minimum size=1.28cm] at (-3.45,0.35) {Present};
  \node[circle, fill=LessonOneSoft, draw=DarkGray, minimum size=1.28cm] at (-1.15,0.35) {Quality};
  \node[circle, fill=LessonThreeSoft,  draw=DarkGray, minimum size=1.35cm] at (1.15,0.35) {Missing\\Pieces};
  \node[circle, fill=CriticalitySoft,  draw=DarkGray, minimum size=1.35cm] at (3.45,0.35) {Future\\Fit};

  \draw[very thick, -{Latex[length=3mm]}, DarkGray] (-3.6,-0.9) -- (3.6,-0.9);
  \node[font=\bfseries, fill=white, inner sep=2pt] at (0,-1.23)
  {Transition planning};
\end{tikzpicture}
\caption{Review frame for evaluating the transition section of the sample IEP.}
\end{figure}

\section*{Completed Indicator 13 Checklist}

\renewcommand{\arraystretch}{1.15}
\small
\begin{longtable}{p{1.65in} p{0.55in} p{2.45in} p{0.8in} p{2.1in}}
\toprule
\rowcolor{SoftBlue}
\textbf{Checklist Item} & \textbf{Present} & \textbf{IEP Evidence} & \textbf{Quality} & \textbf{Brief Evaluation} \\
\midrule
\endfirsthead
\toprule
\rowcolor{SoftBlue}
\textbf{Checklist Item} & \textbf{Present} & \textbf{IEP Evidence} & \textbf{Quality} & \textbf{Brief Evaluation} \\
\midrule
\endhead
Education/training postsecondary goal & Yes & The IEP states that after high school Jimmy will attend college or university to take courses focused on Early Childhood Education. & Moderate & The goal is written as a post-school outcome, but the IEP does not show a clear recent transition assessment explaining why this pathway is the best fit or what supports will be needed. \\
\addlinespace
Employment postsecondary goal & Yes & The IEP states that Jimmy's goal is to be competitively employed working in a day care facility. & Moderate & The goal is observable and future-oriented, but it lacks a developmental pathway such as work-based learning, job shadowing, job coaching, or childcare-related practicum planning. \\
\addlinespace
Independent-living postsecondary goal, when appropriate & Yes & Jimmy stated that his goal is to live in Chicago with friends; his mother shared that he would need some level of support throughout his lifetime. & Low & Independent living is included, but it reads more like a preference statement than a measurable adult-living outcome with support level, safety routines, transportation, budgeting, or health-management expectations. \\
\addlinespace
Postsecondary goals updated annually & Partial & The 2020--2021 IEP includes current postsecondary goals for a 17-year-old student. & Low & Because only one IEP year is available, annual updating cannot be verified across years. \\
\addlinespace
Postsecondary goals based on age-appropriate transition assessment & Partial & The IEP includes academic, cognitive, language, social, parent, health, and functional information; it also states that a functional vocational evaluation is not needed at this time. & Low & The document contains rich evaluation data, but it does not identify a named, recent age-appropriate transition assessment, student interview, vocational interest inventory, adaptive-behavior assessment, or situational vocational assessment tied directly to adult goals. \\
\addlinespace
Transition services/activities reasonably enabling postsecondary goals & Yes & The coordinated activities include instruction, related services, community experiences, employment/adult-living objectives, daily living skills, and a functional vocational assessment decision. & Moderate & Required areas are represented, but several activities are broad and not individualized enough for Jimmy's health-management, self-advocacy, and adult-system navigation needs. \\
\addlinespace
Course of study aligned to goals & Yes & The IEP describes 15:1 math and science, consultant support for English and social studies, Regents Diploma coursework, and supports in core subjects. & Moderate & A course of study is present and connected to diploma progress, but it is not sequenced across remaining school years and does not explicitly include early-childhood, childcare, community-based, or work-based learning coursework. \\
\addlinespace
Annual IEP goals related to transition needs & Partial & Annual goals address reading comprehension, writing, mathematics, and speech/language conversation skills. & Moderate & These goals support school readiness, but they do not directly address self-advocacy, disability disclosure, transportation, health self-management, supported decision-making, or independent living. \\
\addlinespace
Evidence student was invited to IEP meeting where transition was discussed & No & The listed meeting attendees do not include Jimmy, and no invitation documentation appears in the sample. & Low & Student voice appears indirectly in the living goal, but the required evidence of invitation is not visible. \\
\addlinespace
Evidence participating agency was invited with consent, when appropriate & Partial & Transition activities assign responsibility to the school district, Jimmy, and family; no outside agency is named. & Low & The IEP does not document agency invitation or explain why agencies were not invited. For a 17-year-old with complex health and adult support needs, this should be addressed. \\
\bottomrule
\end{longtable}
\normalsize

\section*{What Was Present}

The IEP includes the major transition components expected for a student who is 15 or older. It includes measurable postsecondary goals in education/training, employment, and independent living. Jimmy wants to attend college or university to take courses in Early Childhood Education, become competitively employed in a daycare facility, and live in Chicago with friends. The IEP also includes transition needs, a course of study, coordinated transition activities, annual goals, special education programs and services, accommodations, testing supports, and related services.

The IEP is strongest in its present levels. It documents Jimmy's academic performance, communication profile, social development, parent concerns, health needs, management needs, accommodations, and strengths. It identifies math as a relative strength, explains that Jimmy is friendly and wants to be part of the school community, and describes the support he needs to safely navigate food, temperature regulation, fatigue, communication, and social interactions. That level of detail matters because transition planning cannot be meaningful if it is disconnected from the whole student.

The transition section also includes coordinated activities across several categories. Jimmy will receive instruction in memorization strategies, speech/language services to improve communication, community experience through appointment-making practice, career exploration related to Early Childhood Education, budgeting instruction, and a statement that a functional vocational evaluation is not needed at this time. These items show that the IEP team recognized transition as more than academics alone.

\section*{Quality of the Information Present}

The quality of the transition section is mixed. At a basic compliance level, the IEP is recognizable as an Indicator 13 transition plan because it includes the required categories. However, at a high-quality individualized level, the plan is not yet strong enough for a student with Jimmy's complex profile. The IEP gives a clear picture of school-based support, but it does not fully translate that support into adult-life readiness.

The education and employment goals are meaningful because they reflect Jimmy's stated interest in Early Childhood Education and daycare work. However, the IEP does not show why this pathway has been selected through a current, age-appropriate transition assessment. A high-quality transition plan would include student interview data, interest inventories, adaptive behavior information, community-based observations, or situational vocational assessment data. Without that evidence, the goals appear possible and hopeful, but not yet strongly assessment-driven.

The independent-living goal is the weakest transition goal. Jimmy's desire to live in Chicago with friends is valuable because it gives direct insight into his preferences and sense of future identity. At the same time, the goal does not define the supports he would need to live safely and with dignity. For a chronically ill student with significant food-safety, temperature-regulation, choking-risk, fatigue, and supervision needs, independent living cannot remain vague. The plan should describe likely support levels, health routines, community access, transportation, budgeting, emergency communication, and supported decision-making.

The annual goals also reveal a gap between school success and adult readiness. Reading, writing, math, and speech/language goals support Jimmy's academic progress, but the IEP itself identifies self-advocacy as a transition need. That need does not become a measurable annual goal. For a student who must manage communication access, medical routines, food safety, anxiety, and adult systems, self-advocacy should be central rather than secondary.

\section*{Course Lens: Neurodivergence, Chronic Illness, Depression, and Trauma-Informed Planning}

This review is shaped by the course theme that inclusive education must address more than placement. Students who are neurodivergent, chronically ill, trauma-experienced, depressed, anxious, culturally and linguistically diverse, or navigating multiple identities often need transition plans that protect both access and agency. Support should not become over-direction, and compliance should not replace self-determination.

Jimmy's IEP shows why this matters. His school day requires careful adult support because of serious health needs, food access concerns, temperature-regulation issues, language delays, social vulnerability, and anxiety-related concerns. The plan protects him, and that is necessary. But if transition planning stops at protection, it may not prepare him for adult life. A strong plan would gradually teach Jimmy how to recognize his needs, communicate them, participate in decisions, and use supports without being reduced to those supports.

This is also where depression-responsive and trauma-informed planning become relevant, even when a student does not have a documented depression or trauma diagnosis. Chronic illness, constant supervision, repeated medical vigilance, and social difference can create stress and vulnerability. A transition plan should therefore ask how the student experiences support emotionally. Does the student feel known, or only monitored? Does the plan increase belonging, or does it isolate the student? Does adult support build self-determination, or does it quietly replace the student's decision-making opportunities?

\section*{Additional Information Needed}

The transition section would be stronger if it included a current age-appropriate transition assessment. This should include direct student interview data, family input, interest inventories, adaptive behavior or daily-living assessment, community-based observations, and a situational vocational assessment or explanation of why one is unnecessary. The current IEP has rich present-level information, but transition assessment should directly connect Jimmy's strengths, preferences, interests, support needs, and adult goals.

The IEP should also include clearer student participation documentation. Jimmy's voice appears in the independent-living goal, but there is no evidence that he was invited to the IEP meeting. Indicator 13 transition planning should move toward student-led or at least student-centered planning. Jimmy should have structured opportunities to explain what he wants, what worries him, what supports help him, and what kind of adult life he imagines.

Additional information is also needed around adult systems. Because Jimmy is 17, the team should consider whether vocational rehabilitation, developmental disability services, benefits counseling, supported employment, supported living services, transportation supports, college disability services, or medical transition planning should be introduced. The IEP names school, Jimmy, and family as responsible parties, but it does not show adult-agency coordination.

Finally, the IEP needs more explicit health self-management and self-advocacy planning. Jimmy's adult success will depend on more than academics. It will depend on his ability to communicate health needs, follow safe routines, ask for clarification, manage fatigue, understand food boundaries, identify trusted adults, and use supports in community and postsecondary settings.

\section*{Appropriateness of the Transition Section}

Overall, the transition section is appropriate at a basic compliance level but only partially appropriate at a high-quality individualized level. It includes the required transition categories and connects Jimmy's current school program to broad adult goals. It also reflects care, family involvement, and detailed knowledge of his disability-related needs.

However, the transition plan remains too school-centered and too academically narrow for the complexity of Jimmy's adult-life needs. His goals in college, daycare work, and independent living require more than reading, writing, math, and speech/language progress. They require self-advocacy, disability disclosure, health-management routines, supported decision-making, transportation planning, workplace-readiness practice, and adult-agency coordination.

A stronger plan would not lower expectations. It would make the pathway to those expectations more realistic and more humane. Jimmy's interest in Early Childhood Education should be honored, but it should be supported through structured career exploration, childcare-related observation, communication practice, and careful consideration of health and supervision needs in childcare environments. His goal of living with friends should also be honored, but it should become a measurable adult-living plan rather than remaining only a stated preference.

\section*{Recommended Revisions}

\begin{itemize}
  \item Add a named, dated, age-appropriate transition assessment process, including student interview, interest inventory, adaptive behavior data, and career exploration evidence.
  \item Document that Jimmy was invited to the IEP meeting and provide structured student participation in future transition meetings.
  \item Rewrite the independent-living goal to specify support level, housing expectations, health-management routines, transportation, budgeting, and emergency communication.
  \item Add a measurable annual self-advocacy goal connected to asking for clarification, requesting supports, communicating health needs, and participating in transition decisions.
  \item Reconsider whether a functional vocational evaluation is truly unnecessary, especially given Jimmy's goal of competitive employment in a daycare setting.
  \item Invite or document consideration of adult agencies, including vocational rehabilitation, developmental disability services, college disability supports, benefits counseling, and supported living resources.
  \item Add health-transition planning that gradually builds Jimmy's role in understanding and communicating food, temperature, fatigue, and safety needs.
\end{itemize}

\section*{Sample Improved Goal Language}

\subsection*{Improved Postsecondary Education/Training Goal}

Within one year after exiting high school, Jimmy will enroll in a community college certificate program, supported postsecondary program, or other credentialed training pathway related to Early Childhood Education, and he will register with disability support services before classes begin.

\subsection*{Improved Self-Advocacy Annual Goal}

Given a visual prompt card and guided practice, Jimmy will identify the support he needs in a school, community, or work-preparation task and communicate that need to an adult using spoken language, writing, or a scripted prompt in four out of five opportunities across eight consecutive weeks, as measured by teacher, counselor, or transition-staff logs.

\subsection*{Improved Adult-Living Goal}

Given direct instruction and guided practice, Jimmy will use a step-by-step checklist to schedule an appointment, record the date and materials needed, identify any health or transportation supports required, and follow through with the plan in four out of five trials across school and community-based practice settings.

\section*{Conclusion}

The sample IEP shows a team that knows Jimmy well and takes his health, safety, academic progress, and school participation seriously. That is a strength. The transition section contains the major Indicator 13 elements and provides a basic bridge from school to postsecondary goals. However, the quality of the plan would improve if it moved beyond documenting support and toward building agency.

For neurodivergent, chronically ill, depressed, anxious, trauma-experienced, and otherwise vulnerable students, transition planning should not become paperwork that describes adult life from a distance. It should be a living design for access, dignity, belonging, and self-determination. Jimmy's plan is a solid start, but it needs stronger transition assessment, student voice, adult-agency coordination, independent-living detail, and measurable self-advocacy goals to become truly future-ready.
